%==============================================================================
% MATH 5080 — Introduction to Analysis I (Fall 2026)
% Class:        4
% Date:         Sept 3rd
%==============================================================================

\begin{definition}[Closure]
    The \textbf{closure} of $S$, denoted $\overline{S}$, is defined as 
    \[
    \overline{S} := S \cup \operatorname{bd}  S
    \]
\end{definition}

\begin{definition}[Openess and closedness]
    \hfill
    \begin{enumerate}
        \item $S \subset X$ is called \textbf{open} if $S = \operatorname{int} S$
        \item $S \subset X$ is called \textbf{closed} if $S^c$ is open.
    \end{enumerate}
\end{definition}

\begin{remark}
    A point cannot be both an interior point and a boundary point.
\end{remark}

\begin{remark}
    A set can be both open and closed. For example 
    \begin{itemize}
        \item $\mathbb{R}$ is both open and closed.
        \item $\emptyset$ is both open and closed.
        \item $\mathbb{N}$ is not open, but closed.
    \end{itemize} 
\end{remark}

\begin{example}
    $S = (1, 3)$ is open and not closed.
\end{example}

\begin{example}
    Let $(X, \sigma)$ be a metric space where $\sigma$ is the discrete 
    metric. Take $x_0 \in X$, $S = \{ x_0 \} $ is open. 
    \medskip 

    To see why, let $\ep = \frac{1}{2}$. 
    \begin{align*}
        & N \left(x_0, \frac{1}{2} \right) = \{ x_0 \}  \subset S \\
        \implies & x_0 \in \operatorname{int}  S \\
        \implies & \operatorname{int}  S = S \\
        \implies & S \text{ open}
    \end{align*}
    Hence any point in a set is interior, and any set is open and closed.
\end{example}

\begin{proposition}
    \hfill
    \begin{enumerate}
        \item $S$ open $\iff$ $\operatorname{bd} S \subset S^c$
        \item $S$ closed 
$\iff \operatorname{bd} S \subset S$
$\iff S = \overline{S} := S \cup \operatorname{bd} S$ 
        \item $\overline{S}$ is closed
        \item $\operatorname{int} S$ is open $\implies \operatorname{int} \left( \operatorname{int} (S)  \right) = \operatorname{int}(S)  $
    \end{enumerate}
\end{proposition}

\begin{proof}[Proof. $S$ closed $\iff$ $\operatorname{bd} S \subset S$]
    \hfill

    By definition, $S$ closed $\iff$ $S^c$ open. 
    Also by definition, $\operatorname{bd} S = \operatorname{bd} S^c$.
    \begin{align*}
        \operatorname{bd} S \subset S & \iff \operatorname{bd} S^c \subset S && \text{ by definition} \\
        & \iff \operatorname{bd} S^c \cap S^c = \emptyset \\
        & \iff \forall x \in S^c, x \notin \operatorname{bd} S^c \\
        & \iff \forall x \in S^c, x \in \operatorname{int} S^c \\
        & \iff S^c = \operatorname{int} S^c \\
        & \iff S^c \text{ open} \\
        & \iff S \text{ closed}
    \end{align*}
\end{proof}

\begin{proposition}
    \hfill
    \begin{enumerate}
        \item The union of open sets is open.  
        \item The intersection of finitely many open sets is open.
    \end{enumerate}
    \remind{Add in analogous propositions about closed sets.}
\end{proposition}

\begin{example}
    $S_n = \left( - \frac{1}{n}, \frac{1}{n} \right) $ is open for all $n \in \mathbb{N}$. 
    $\bigcap_{n=1}^{\infty} S_n = \{ 0 \} $ is closed. 

    \medskip 

    To see why the intersection is closed, 
    $\forall x >0$, by AP, $\exists n \in \mathbb{N}$ 
    s.t. $0 < \frac{1}{n} < x$. 
    Hence $x \notin S_n$ and 
    $S \notin \bigcap S_n$. 
    \medskip

    At the same time, $0 \in S_n, \forall n \in \mathbb{N}$. 
    \medskip

    Since $0 \notin \operatorname{int} \{ 0 \}  $, $\{ 0 \} $ is not open.
\end{example}

\begin{proof}[Proof of 3.9.2]
    WTS intersection of finitely many open sets is open. 
    \medskip

    Let $S_1, \hdots S_n$ be open sets, $S = \bigcap_{k = 1}^n S_k$. 
    WTS $\forall x S, \exists \ep \in \mathbb{R}, \ep > 0$
    s.t. 
    $N(x , \ep) \subset S$. 
    Note that 
    \[
        N(x, \ep ) \subset S \iff
        N(x, \ep ) \subset S_i, \forall i = 1, \hdots, n
    \]
    Since $x \in S_i$ and $S_i$ is open, $\exists \ep_i > 0$ s.t. 
    $N(x, \ep_i) \subset S_i$. 
    Let $\ep = \min \{ \ep_1, \hdots, \ep_n  \} $,
    and we know such $\ep$ exists 
    because a sequence of finitely many numbers has a minimum.
\end{proof}

\begin{definition}[Accumulation points and isolated points]
    \hfill 
    \begin{enumerate}
        \item
    Let $S \subset X$, a point $x \in X$ is an \textbf{accumulation point}
    of $S$ if $\forall \ep > 0$, 
    \[
        N^*(x, \ep) \cap S \neq \emptyset
    \]
    The set of all accumulation points of $S$ is denoted $S^{\prime}$.
        \item If $x \in S \setminus S^{\prime}$, then $x$ is an 
\textbf{isolated point} of $S$.
    \end{enumerate}
\end{definition}

\begin{example}
    $S = \{ 0 \}  \cup (1, 2)$, $S^{\prime} = [1, 2]$
\end{example}

\begin{example}
    $S = \mathbb{Q} \cap (1, 2)$, $S^{\prime} = [1, 2]$
\end{example}

\begin{proposition}\label{prop:closure-acc}
    \hfill
    \begin{enumerate}[ref=\theproposition.\arabic*]
        \item\label{prop:closure-acc-1} $\overline{S} = S \cup S^{\prime}$
        \item\label{prop:closure-acc-2} $S$ close $\iff S^{\prime} \subset S$
        \item\label{prop:closure-acc-3} $\overline{S}$ is the smallest closed set that contains $S$.
    \end{enumerate}
\end{proposition}

\begin{proof}[Proof of \ref{prop:closure-acc-1}]
    WTS $S \cup \operatorname{bd} S = S \cup S^{\prime}$. 
    \begin{enumerate}
        \item We first show that $LHS \subset RHS$. 
$\forall x \in S \cup \operatorname{bd} S$, $x \in S \cup S^{\prime}$. 
\medskip 

If $x \in S$, done. 
\medskip

If $x \in \operatorname{bd} S$, $x \notin S$ then $x \in \operatorname{bd} S \cap S^c$. 
WTS $x \in S^{\prime}$, i.e. $\forall \ep > 0, N^*(x, \ep) \neq \emptyset$. 
\medskip 

Since $x \in \operatorname{bd} S$, we know $N(x, \ep) \cap S \neq \emptyset$. 
Since $x \in S^c$, we know $\exists y \in N(x, \ep)\cap S, y \neq x$. 
Hence $y \in N^*(x, \ep) \cap S$, 
and $N^*(x, \ep) \neq \emptyset$.
\item We then show $RHS \subset LHS$, i.e. 
$\forall x \in S \cup S^{\prime}, x \in S \cup \operatorname{bd} S$. 
\medskip 

If $x \in S^{\prime}$, done. 
\medskip 

If $x \in S^{\prime}, x \notin S$, $x \in S^{\prime} \cap S^c$, 
WTS $x \in \operatorname{bd} S$
\remind{Exercise}.
    \end{enumerate}
\end{proof}











